\documentclass{article}
\usepackage{amsmath, amssymb, geometry, enumitem}
\geometry{a4paper, margin=1in}

\title{Math Problem Sets: The Inverse of a Matrix}
\author{}
\date{}

\begin{document}

\maketitle

\tableofcontents
\newpage

\section{Concept Checklist \& Question Mapping}

To ensure complete coverage of the topics derived from the provided reading, the problems are categorized into 10 fundamental mathematical concepts. The distribution heavily weighs advanced and conceptual techniques to reinforce deep understanding.

\begin{itemize}
    \item \textbf{Concept 1: $2 \times 2$ Inverse Formula Application.} Directly applying the definition $1/(ad-bc) \begin{bmatrix} d & -b \\ -c & a \end{bmatrix}$ and checking if the matrix is invertible. \\
    \textit{Relevant Questions: 1, 2, 3, 4, 31}
    
    \item \textbf{Concept 2: Verification and Definition Testing.} Utilizing the baseline definition $AA^{-1} = I$ to prove matrices are inverses without recalculating them. \\
    \textit{Relevant Questions: 5, 6}
    
    \item \textbf{Concept 3: Solving Linear Systems via Inverses.} Translating systems to $Ax = b$ and exploiting pre-calculated inverses to find $x = A^{-1}b$. \\
    \textit{Relevant Questions: 7, 8, 9}
    
    \item \textbf{Concept 4: $n \times n$ Inverse Algorithm.} Utilizing standard Gaussian Elimination on the augmented matrix $[A \mid I]$ to yield $[I \mid A^{-1}]$. \\
    \textit{Relevant Questions: 10, 11, 12}
    
    \item \textbf{Concept 5: Targeted Inverse Columns.} Augmenting $A$ solely with a standard basis vector $e_j$ to retrieve a specific column of the inverse, saving computational resources. \\
    \textit{Relevant Questions: 13, 14}
    
    \item \textbf{Concept 6: Matrix Equations and Reversals.} Isolating unknown matrices using correct left or right multiplication (e.g., $X = B A^{-1}$). \\
    \textit{Relevant Questions: 15, 16}
    
    \item \textbf{Concept 7: Algebraic Proofs.} Symbolic manipulation utilizing inverse properties, un-distributing terms, and rigorously validating steps using matrix axioms. \\
    \textit{Relevant Questions: 17, 18, 19, 32}
    
    \item \textbf{Concept 8: Conceptual Synthesis (True/False).} Diagnostic checks on theorems, null spaces, and non-linear properties of inverses. \\
    \textit{Relevant Questions: 20, 21, 22, 23, 24}
    
    \item \textbf{Concept 9: Block Matrix Inversion.} Advanced technique of reciprocating individual square blocks across the main diagonal. \\
    \textit{Relevant Questions: 25, 26}
    
    \item \textbf{Concept 10: Find the Flaw (Diagnostic Testing).} Identifying, explaining, and correcting common algebraic illusions, order-of-operation violations, and notation fallacies. \\
    \textit{Relevant Questions: 27, 28, 29, 30}
\end{itemize}

\newpage

\section{Problem Sets}

\textbf{Direct Formula Application} \\
\textbf{1.} Find the inverse of the matrix $A = \begin{bmatrix} 7 & 4 \\ 5 & 3 \end{bmatrix}$, if it exists. \\
\textbf{2.} Find the inverse of the matrix $B = \begin{bmatrix} -2 & 5 \\ 3 & -8 \end{bmatrix}$, if it exists. \\
\textbf{3.} Find the inverse of the matrix $C = \begin{bmatrix} 4 & -2 \\ -6 & 3 \end{bmatrix}$, if it exists. \\
\textbf{4.} Find the inverse of the matrix $D = \begin{bmatrix} 6 & 1 \\ 4 & 2 \end{bmatrix}$, if it exists. \\
\textbf{31.} Find the inverse of the matrix $E = \begin{bmatrix} 0 & 2 \\ 3 & 4 \end{bmatrix}$, if it exists.

\vspace{10pt}
\noindent
\textbf{Verification and Definition Testing} \\
\textbf{5.} Verify that $M = \begin{bmatrix} 2 & 1 \\ 5 & 3 \end{bmatrix}$ and $N = \begin{bmatrix} 3 & -1 \\ -5 & 2 \end{bmatrix}$ are inverses by computing their product. \\
\textbf{6.} Verify that the matrix $A = \begin{bmatrix} 4 & 5 \\ 2 & 3 \end{bmatrix}$ has the inverse $A^{-1} = \begin{bmatrix} 1.5 & -2.5 \\ -1 & 2 \end{bmatrix}$ by utilizing the foundational definition of an inverse.

\vspace{10pt}
\noindent
\textbf{Solving Linear Systems via Inverses} \\
\textbf{7.} Use the inverse found in Problem 1 to solve the following linear system:
\begin{align*}
7x_1 + 4x_2 &= 5 \\
5x_1 + 3x_2 &= 2
\end{align*}
\textbf{8.} Use an inverse matrix to solve the following linear system:
\begin{align*}
-2x_1 + 5x_2 &= -1 \\
3x_1 - 8x_2 &= 4
\end{align*}
\textbf{9.} Suppose $A = \begin{bmatrix} 6 & 1 \\ 4 & 2 \end{bmatrix}$. Using the inverse found in Problem 4, simultaneously solve the matrix equation $AX = B$ for the vectors $b_1 = \begin{bmatrix} 8 \\ 0 \end{bmatrix}$ and $b_2 = \begin{bmatrix} 2 \\ 4 \end{bmatrix}$. 

\vspace{10pt}
\noindent
\textbf{Algorithm for Inverting $n \times n$ Matrices} \\
\textbf{10.} Find the inverse of $A = \begin{bmatrix} 1 & -1 & 0 \\ 2 & 0 & 1 \\ 0 & 2 & 2 \end{bmatrix}$ using the augmented matrix algorithm $[A \mid I]$. \\
\textbf{11.} Find the inverse of $B = \begin{bmatrix} 2 & 3 & 4 \\ 2 & 1 & 1 \\ -1 & 1 & 2 \end{bmatrix}$ using the augmented matrix algorithm $[B \mid I]$. \\
\textbf{12.} Find the inverse of $C = \begin{bmatrix} 1 & 0 & 2 \\ 0 & 3 & 0 \\ 4 & 0 & 5 \end{bmatrix}$ using the augmented matrix algorithm. 

\vspace{10pt}
\noindent
\textbf{Targeted Inverse Columns} \\
\textbf{13.} Let $A = \begin{bmatrix} 2 & 1 & 3 \\ 1 & -1 & 1 \\ 3 & 2 & 5 \end{bmatrix}$. Find the second column of $A^{-1}$ without computing the rest of the inverse. \\
\textbf{14.} Let $M = \begin{bmatrix} 1 & 2 & -1 \\ -2 & -3 & 1 \\ 1 & 4 & -4 \end{bmatrix}$. Find the third column of $M^{-1}$ without computing the rest of the inverse. 

\vspace{10pt}
\noindent
\textbf{Matrix Equations and Reversals} \\
\textbf{15.} Solve the matrix equation $XA = B$ for the unknown matrix $X$, assuming $A = \begin{bmatrix} 3 & 5 \\ 1 & 2 \end{bmatrix}$ and $B = \begin{bmatrix} 4 & -1 \\ 2 & 0 \end{bmatrix}$. \\
\textbf{16.} Solve the matrix equation $AY + B = C$ for the unknown matrix $Y$, where $A = \begin{bmatrix} 1 & -2 \\ -1 & 3 \end{bmatrix}$, $B = \begin{bmatrix} 1 & 2 \\ 3 & 4 \end{bmatrix}$, and $C = \begin{bmatrix} 5 & 6 \\ 7 & 8 \end{bmatrix}$. 

\vspace{10pt}
\noindent
\textbf{Algebraic Proofs} \\
\textbf{17.} Prove that if $P$ and $Q$ are invertible $n \times n$ matrices, then $(P^T Q)^{-1} = Q^{-1} (P^{-1})^T$. \\
\textbf{18.} Use matrix algebra to show that if $A$ is an invertible matrix and satisfies the polynomial equation $A^2 - 3A + 2I = 0$, then $A^{-1} = \frac{3}{2}I - \frac{1}{2}A$. \\
\textbf{19.} Assume $A, B$, and $X$ are $n \times n$ matrices. Furthermore, assume $A, B$, and $(A - AX)$ are invertible. Simplify the expression for the matrix $X$ if $(A - AX)^{-1} = B^{-1}$. \\
\textbf{32.} Let $A$ be an invertible $n \times n$ matrix. Prove that $(A^T)^{-1} = (A^{-1})^T$ relying strictly on matrix axioms. 

\vspace{10pt}
\noindent
\textbf{Conceptual Synthesis (True/False)} \\
\textit{Determine if the statement is True or False. Briefly justify your answer.} \\
\textbf{20.} (T/F) If $A$ and $B$ are $n \times n$ matrices and their product $AB$ is invertible, then $A$ and $B$ must both individually be invertible. \\
\textbf{21.} (T/F) If $A$ is an invertible $n \times n$ matrix, then the homogeneous system $Ax = 0$ has only the trivial solution $x = 0$. \\
\textbf{22.} (T/F) The inverse of the sum of two invertible matrices $A$ and $B$ is identically the sum of their inverses: $(A+B)^{-1} = A^{-1} + B^{-1}$. \\
\textbf{23.} (T/F) If $A$ is a singular (non-invertible) $n \times n$ matrix, then its row-reduced echelon form contains at least one row of all zeros. \\
\textbf{24.} (T/F) For any non-zero scalar $k$ and invertible matrix $A$, the inverse evaluates to $(kA)^{-1} = k A^{-1}$.

\vspace{10pt}
\noindent
\textbf{Block Matrix Inversion} \\
\textbf{25.} Find the inverse of the $4 \times 4$ block diagonal matrix $M = \begin{bmatrix} A & 0 \\ 0 & B \end{bmatrix}$ where $A = \begin{bmatrix} 2 & 1 \\ 5 & 3 \end{bmatrix}$ and $B = \begin{bmatrix} 4 & 0 \\ 0 & 5 \end{bmatrix}$, and $0$ represents $2 \times 2$ zero matrices. \\
\textbf{26.} Find the inverse of the matrix $P = \begin{bmatrix} 1 & 0 & 0 & 0 \\ 0 & 1 & 0 & 0 \\ 0 & 0 & 3 & 4 \\ 0 & 0 & 2 & 3 \end{bmatrix}$. 

\vspace{10pt}
\noindent
\textbf{Diagnostic Testing: "Find the Flaw"} \\
\textit{Each of the following solutions contains a single, subtle conceptual or algebraic error. Spot the mistake, explain why it is mathematically invalid, and provide the correct solution.} \\
\textbf{27.} \textbf{Question:} Simplify $(ABC)^{-1}$, assuming $A, B$, and $C$ are invertible $n \times n$ matrices. \\
\textbf{Flawed Solution:} By the property of inverses, we apply the inverse to each matrix in reverse order: $(ABC)^{-1} = A^{-1} B^{-1} C^{-1}$. 

\vspace{5pt}
\noindent
\textbf{28.} \textbf{Question:} Solve the matrix equation $AX = B$ for the unknown matrix $X$, assuming $A$ is invertible. \\
\textbf{Flawed Solution:} To isolate $X$, we must divide both sides by the matrix $A$. Therefore, $X = \frac{B}{A}$.

\vspace{5pt}
\noindent
\textbf{29.} \textbf{Question:} Find the inverse of $A = \begin{bmatrix} 7 & 2 \\ 3 & 1 \end{bmatrix}$. \\
\textbf{Flawed Solution:} The determinant is $(7)(1) - (2)(3) = 1$. The inverse is found by swapping all elements and negating the main diagonal: $A^{-1} = \begin{bmatrix} -1 & 3 \\ 2 & -7 \end{bmatrix}$. 

\vspace{5pt}
\noindent
\textbf{30.} \textbf{Question:} Find the inverse of the matrix $M = \frac{1}{2} A$, given that $A^{-1} = \begin{bmatrix} 4 & 2 \\ 6 & 8 \end{bmatrix}$. \\
\textbf{Flawed Solution:} The scalar factors out directly when dealing with inverses. $M^{-1} = \frac{1}{2} A^{-1} = \begin{bmatrix} 2 & 1 \\ 3 & 4 \end{bmatrix}$.

\newpage

\section{Detailed Solutions}

\textbf{1.} $det(A) = (7)(3) - (4)(5) = 21 - 20 = 1 \neq 0$. \\
$A^{-1} = \frac{1}{1} \begin{bmatrix} 3 & -4 \\ -5 & 7 \end{bmatrix} = \begin{bmatrix} 3 & -4 \\ -5 & 7 \end{bmatrix}$.

\vspace{5pt}
\noindent
\textbf{2.} $det(B) = (-2)(-8) - (5)(3) = 16 - 15 = 1 \neq 0$. \\
$B^{-1} = \frac{1}{1} \begin{bmatrix} -8 & -5 \\ -3 & -2 \end{bmatrix} = \begin{bmatrix} -8 & -5 \\ -3 & -2 \end{bmatrix}$.

\vspace{5pt}
\noindent
\textbf{3.} $det(C) = (4)(3) - (-2)(-6) = 12 - 12 = 0$. \\
Because the determinant is zero, $C$ is a singular matrix and \textbf{is not invertible}.

\vspace{5pt}
\noindent
\textbf{4.} $det(D) = (6)(2) - (1)(4) = 12 - 4 = 8 \neq 0$. \\
$D^{-1} = \frac{1}{8} \begin{bmatrix} 2 & -1 \\ -4 & 6 \end{bmatrix} = \begin{bmatrix} 1/4 & -1/8 \\ -1/2 & 3/4 \end{bmatrix}$.

\vspace{5pt}
\noindent
\textbf{31.} $det(E) = (0)(4) - (2)(3) = -6 \neq 0$. \\
$E^{-1} = -\frac{1}{6} \begin{bmatrix} 4 & -2 \\ -3 & 0 \end{bmatrix} = \begin{bmatrix} -2/3 & 1/3 \\ 1/2 & 0 \end{bmatrix}$.

\vspace{5pt}
\noindent
\textbf{5.} $MN = \begin{bmatrix} 2(3)+1(-5) & 2(-1)+1(2) \\ 5(3)+3(-5) & 5(-1)+3(2) \end{bmatrix} = \begin{bmatrix} 6-5 & -2+2 \\ 15-15 & -5+6 \end{bmatrix} = \begin{bmatrix} 1 & 0 \\ 0 & 1 \end{bmatrix}$. \\
Since $MN = I$, they are verified as inverses.

\vspace{5pt}
\noindent
\textbf{6.} $A^{-1}A = \begin{bmatrix} 1.5(4) - 2.5(2) & 1.5(5) - 2.5(3) \\ -1(4) + 2(2) & -1(5) + 2(3) \end{bmatrix} = \begin{bmatrix} 6-5 & 7.5-7.5 \\ -4+4 & -5+6 \end{bmatrix} = \begin{bmatrix} 1 & 0 \\ 0 & 1 \end{bmatrix}$.

\vspace{5pt}
\noindent
\textbf{7.} Formulate $Ax = b \implies x = A^{-1}b$. From Problem 1, $A^{-1} = \begin{bmatrix} 3 & -4 \\ -5 & 7 \end{bmatrix}$. \\
$x = \begin{bmatrix} 3 & -4 \\ -5 & 7 \end{bmatrix} \begin{bmatrix} 5 \\ 2 \end{bmatrix} = \begin{bmatrix} 15 - 8 \\ -25 + 14 \end{bmatrix} = \begin{bmatrix} 7 \\ -11 \end{bmatrix}$. Final Answer: $x_1 = 7, x_2 = -11$.

\vspace{5pt}
\noindent
\textbf{8.} Formulate $Bx = b$. From Problem 2, $B^{-1} = \begin{bmatrix} -8 & -5 \\ -3 & -2 \end{bmatrix}$. \\
$x = \begin{bmatrix} -8 & -5 \\ -3 & -2 \end{bmatrix} \begin{bmatrix} -1 \\ 4 \end{bmatrix} = \begin{bmatrix} 8 - 20 \\ 3 - 8 \end{bmatrix} = \begin{bmatrix} -12 \\ -5 \end{bmatrix}$. Final Answer: $x_1 = -12, x_2 = -5$.

\vspace{5pt}
\noindent
\textbf{9.} Formulate block system $X = D^{-1} \begin{bmatrix} b_1 \mid b_2 \end{bmatrix}$. \\
$X = \frac{1}{8} \begin{bmatrix} 2 & -1 \\ -4 & 6 \end{bmatrix} \begin{bmatrix} 8 & 2 \\ 0 & 4 \end{bmatrix} = \frac{1}{8} \begin{bmatrix} 16 & 0 \\ -32 & 16 \end{bmatrix} = \begin{bmatrix} 2 & 0 \\ -4 & 2 \end{bmatrix}$. \\
Solution 1 is $(2, -4)$, Solution 2 is $(0, 2)$.

\vspace{5pt}
\noindent
\textbf{10.} Augment $[A \mid I]$:
$\begin{bmatrix} 1 & -1 & 0 & \mid & 1 & 0 & 0 \\ 2 & 0 & 1 & \mid & 0 & 1 & 0 \\ 0 & 2 & 2 & \mid & 0 & 0 & 1 \end{bmatrix}$ \\
$R_2 \to R_2 - 2R_1 \implies \begin{bmatrix} 1 & -1 & 0 & \mid & 1 & 0 & 0 \\ 0 & 2 & 1 & \mid & -2 & 1 & 0 \\ 0 & 2 & 2 & \mid & 0 & 0 & 1 \end{bmatrix}$ \\
$R_3 \to R_3 - R_2 \implies \begin{bmatrix} 1 & -1 & 0 & \mid & 1 & 0 & 0 \\ 0 & 2 & 1 & \mid & -2 & 1 & 0 \\ 0 & 0 & 1 & \mid & 2 & -1 & 1 \end{bmatrix}$ \\
$R_2 \to R_2 - R_3 \implies \begin{bmatrix} 1 & -1 & 0 & \mid & 1 & 0 & 0 \\ 0 & 2 & 0 & \mid & -4 & 2 & -1 \\ 0 & 0 & 1 & \mid & 2 & -1 & 1 \end{bmatrix}$ \\
$R_2 \to \frac{1}{2}R_2 \implies \begin{bmatrix} 1 & -1 & 0 & \mid & 1 & 0 & 0 \\ 0 & 1 & 0 & \mid & -2 & 1 & -0.5 \\ 0 & 0 & 1 & \mid & 2 & -1 & 1 \end{bmatrix}$ \\
$R_1 \to R_1 + R_2 \implies \begin{bmatrix} 1 & 0 & 0 & \mid & -1 & 1 & -0.5 \\ 0 & 1 & 0 & \mid & -2 & 1 & -0.5 \\ 0 & 0 & 1 & \mid & 2 & -1 & 1 \end{bmatrix}$. \\
Final Answer: $A^{-1} = \begin{bmatrix} -1 & 1 & -0.5 \\ -2 & 1 & -0.5 \\ 2 & -1 & 1 \end{bmatrix}$.

\vspace{5pt}
\noindent
\textbf{11.} Augment $[B \mid I]$:
$\begin{bmatrix} 2 & 3 & 4 & \mid & 1 & 0 & 0 \\ 2 & 1 & 1 & \mid & 0 & 1 & 0 \\ -1 & 1 & 2 & \mid & 0 & 0 & 1 \end{bmatrix}$ \\
Perform standard Gaussian Elimination (Row swap, scaling, reduction): \\
$R_1 \leftrightarrow R_3, R_1 \to -R_1 \implies \begin{bmatrix} 1 & -1 & -2 & \mid & 0 & 0 & -1 \\ 2 & 1 & 1 & \mid & 0 & 1 & 0 \\ 2 & 3 & 4 & \mid & 1 & 0 & 0 \end{bmatrix}$ \\
$R_2 \to R_2 - 2R_1, R_3 \to R_3 - 2R_1 \implies \begin{bmatrix} 1 & -1 & -2 & \mid & 0 & 0 & -1 \\ 0 & 3 & 5 & \mid & 0 & 1 & 2 \\ 0 & 5 & 8 & \mid & 1 & 0 & 2 \end{bmatrix}$ \\
Upon complete systematic zeroing of above and below diagonals, the right side renders: \\
$B^{-1} = \begin{bmatrix} -1 & 2 & 1 \\ 5 & -8 & -6 \\ -3 & 5 & 4 \end{bmatrix}$.

\vspace{5pt}
\noindent
\textbf{12.} Augment $[C \mid I]$:
$\begin{bmatrix} 1 & 0 & 2 & \mid & 1 & 0 & 0 \\ 0 & 3 & 0 & \mid & 0 & 1 & 0 \\ 4 & 0 & 5 & \mid & 0 & 0 & 1 \end{bmatrix}$ \\
$R_2 \to \frac{1}{3}R_2$ and $R_3 \to R_3 - 4R_1 \implies \begin{bmatrix} 1 & 0 & 2 & \mid & 1 & 0 & 0 \\ 0 & 1 & 0 & \mid & 0 & 1/3 & 0 \\ 0 & 0 & -3 & \mid & -4 & 0 & 1 \end{bmatrix}$ \\
$R_3 \to -\frac{1}{3}R_3 \implies \begin{bmatrix} 1 & 0 & 2 & \mid & 1 & 0 & 0 \\ 0 & 1 & 0 & \mid & 0 & 1/3 & 0 \\ 0 & 0 & 1 & \mid & 4/3 & 0 & -1/3 \end{bmatrix}$ \\
$R_1 \to R_1 - 2R_3 \implies \begin{bmatrix} 1 & 0 & 0 & \mid & -5/3 & 0 & 2/3 \\ 0 & 1 & 0 & \mid & 0 & 1/3 & 0 \\ 0 & 0 & 1 & \mid & 4/3 & 0 & -1/3 \end{bmatrix}$. \\
Final Answer: $C^{-1} = \begin{bmatrix} -5/3 & 0 & 2/3 \\ 0 & 1/3 & 0 \\ 4/3 & 0 & -1/3 \end{bmatrix}$.

\vspace{5pt}
\noindent
\textbf{13.} To find the second column, solve $Ax = e_2 = \begin{bmatrix} 0 \\ 1 \\ 0 \end{bmatrix}$. \\
Augment: $\begin{bmatrix} 2 & 1 & 3 & \mid & 0 \\ 1 & -1 & 1 & \mid & 1 \\ 3 & 2 & 5 & \mid & 0 \end{bmatrix}$. Row reducing yields $x = \begin{bmatrix} -1 \\ -1 \\ 1 \end{bmatrix}$. The second column is $\begin{bmatrix} -1 \\ -1 \\ 1 \end{bmatrix}$.

\vspace{5pt}
\noindent
\textbf{14.} To find the third column, solve $Mx = e_3 = \begin{bmatrix} 0 \\ 0 \\ 1 \end{bmatrix}$. \\
Augment: $\begin{bmatrix} 1 & 2 & -1 & \mid & 0 \\ -2 & -3 & 1 & \mid & 0 \\ 1 & 4 & -4 & \mid & 1 \end{bmatrix}$. Row reducing yields $x_3 = -1, x_2 = -1, x_1 = 1$. The column is $\begin{bmatrix} 1 \\ -1 \\ -1 \end{bmatrix}$.

\vspace{5pt}
\noindent
\textbf{15.} Multiply both sides on the right by $A^{-1}$: $X(AA^{-1}) = BA^{-1} \implies X = BA^{-1}$. \\
$A^{-1} = \frac{1}{6-5} \begin{bmatrix} 2 & -5 \\ -1 & 3 \end{bmatrix} = \begin{bmatrix} 2 & -5 \\ -1 & 3 \end{bmatrix}$. \\
$X = \begin{bmatrix} 4 & -1 \\ 2 & 0 \end{bmatrix} \begin{bmatrix} 2 & -5 \\ -1 & 3 \end{bmatrix} = \begin{bmatrix} (8+1) & (-20-3) \\ (4+0) & (-10+0) \end{bmatrix} = \begin{bmatrix} 9 & -23 \\ 4 & -10 \end{bmatrix}$.

\vspace{5pt}
\noindent
\textbf{16.} Isolate $AY$: $AY = C - B = \begin{bmatrix} 5 & 6 \\ 7 & 8 \end{bmatrix} - \begin{bmatrix} 1 & 2 \\ 3 & 4 \end{bmatrix} = \begin{bmatrix} 4 & 4 \\ 4 & 4 \end{bmatrix}$. \\
Multiply both sides on the left by $A^{-1}$. \\
$A^{-1} = \frac{1}{3-2} \begin{bmatrix} 3 & 2 \\ 1 & 1 \end{bmatrix} = \begin{bmatrix} 3 & 2 \\ 1 & 1 \end{bmatrix}$. \\
$Y = A^{-1} \begin{bmatrix} 4 & 4 \\ 4 & 4 \end{bmatrix} = \begin{bmatrix} 3 & 2 \\ 1 & 1 \end{bmatrix} \begin{bmatrix} 4 & 4 \\ 4 & 4 \end{bmatrix} = \begin{bmatrix} (12+8) & (12+8) \\ (4+4) & (4+4) \end{bmatrix} = \begin{bmatrix} 20 & 20 \\ 8 & 8 \end{bmatrix}$.

\vspace{5pt}
\noindent
\textbf{17.} To prove that $Q^{-1} (P^{-1})^T$ is the inverse of $(P^T Q)$, we multiply them and see if the product is $I$. \\
$(P^T Q) \big(Q^{-1} (P^{-1})^T \big) = P^T (Q Q^{-1}) (P^{-1})^T = P^T (I) (P^{-1})^T = P^T (P^T)^{-1}$ (using the property that the transpose of an inverse is the inverse of a transpose). \\
$= P^T (P^T)^{-1} = I$. The proof is complete.

\vspace{5pt}
\noindent
\textbf{18.} Given $A^2 - 3A + 2I = 0$, isolate the identity matrix: $2I = 3A - A^2$. \\
Left-multiply by $A^{-1}$: $A^{-1}(2I) = A^{-1}(3A - A^2) \implies 2A^{-1} = 3I - A$. \\
Divide by 2: $A^{-1} = \frac{3}{2}I - \frac{1}{2}A$. The proof is complete.

\vspace{5pt}
\noindent
\textbf{19.} Start with $(A - AX)^{-1} = B^{-1}$. Invert both sides: $A - AX = B$. \\
Isolate $AX$: $AX = A - B$. \\
Left-multiply by $A^{-1}$: $X = A^{-1}(A - B)$. \\
Distribute: $X = I - A^{-1}B$. 

\vspace{5pt}
\noindent
\textbf{32.} Prove that $(A^T)^{-1} = (A^{-1})^T$. We must show that $A^T$ multiplied by $(A^{-1})^T$ equals $I$. \\
Using the fundamental transpose rule $(XY)^T = Y^T X^T$: \\
$A^T (A^{-1})^T = (A^{-1} A)^T = (I)^T = I$. The proof is complete.

\vspace{5pt}
\noindent
\textbf{20.} \textbf{True.} By properties of the determinant, $det(AB) = det(A)det(B)$. If $AB$ is invertible, its determinant is non-zero, demanding that $det(A) \neq 0$ and $det(B) \neq 0$, making both individually invertible.

\vspace{5pt}
\noindent
\textbf{21.} \textbf{True.} We left multiply the equation by $A^{-1}$: $A^{-1}Ax = A^{-1}0 \implies Ix = 0 \implies x = 0$.

\vspace{5pt}
\noindent
\textbf{22.} \textbf{False.} The matrix inverse operation does not distribute over addition. This is the "Additive Illusion" trap. Generally, $(A+B)^{-1} \neq A^{-1} + B^{-1}$.

\vspace{5pt}
\noindent
\textbf{23.} \textbf{True.} By the Invertible Matrix Theorem, a square matrix is invertible if and only if it is row equivalent to the identity matrix $I$. If it is singular, it cannot fully reduce to $I$, which mathematically dictates the presence of at least one row of all zeros in its RREF.

\vspace{5pt}
\noindent
\textbf{24.} \textbf{False.} The inverse of a scalar-multiplied matrix requires the reciprocal of the scalar. The correct property is $(kA)^{-1} = \frac{1}{k} A^{-1}$.

\vspace{5pt}
\noindent
\textbf{25.} By block matrix rules, $M^{-1} = \begin{bmatrix} A^{-1} & 0 \\ 0 & B^{-1} \end{bmatrix}$. \\
$A^{-1} = \frac{1}{6-5}\begin{bmatrix} 3 & -1 \\ -5 & 2 \end{bmatrix} = \begin{bmatrix} 3 & -1 \\ -5 & 2 \end{bmatrix}$. $B$ is a diagonal matrix, its inverse is the reciprocal of the diagonals: $B^{-1} = \begin{bmatrix} 1/4 & 0 \\ 0 & 1/5 \end{bmatrix}$. \\
Final Answer: $M^{-1} = \begin{bmatrix} 3 & -1 & 0 & 0 \\ -5 & 2 & 0 & 0 \\ 0 & 0 & 1/4 & 0 \\ 0 & 0 & 0 & 1/5 \end{bmatrix}$.

\vspace{5pt}
\noindent
\textbf{26.} Partition $P$ into a block diagonal matrix where block $A = I_2$ and block $B = \begin{bmatrix} 3 & 4 \\ 2 & 3 \end{bmatrix}$. \\
$A^{-1} = I_2$. $B^{-1} = \frac{1}{9-8} \begin{bmatrix} 3 & -4 \\ -2 & 3 \end{bmatrix} = \begin{bmatrix} 3 & -4 \\ -2 & 3 \end{bmatrix}$. \\
Final Answer: $P^{-1} = \begin{bmatrix} 1 & 0 & 0 & 0 \\ 0 & 1 & 0 & 0 \\ 0 & 0 & 3 & -4 \\ 0 & 0 & -2 & 3 \end{bmatrix}$.

\vspace{5pt}
\noindent
\textbf{27.} \textbf{Flaw:} The student distributed the inverse operator without actually reversing the matrix order (The "Shoes and Socks" Reverse-Order Law). \\
\textbf{Correct Solution:} The inverse of a product is the product of their inverses taken in reverse order. $(ABC)^{-1} = C^{-1} B^{-1} A^{-1}$.

\vspace{5pt}
\noindent
\textbf{28.} \textbf{Flaw:} The student attempted to execute matrix division, an operation that is strictly undefined in linear algebra. Multiplicative cancellation must be executed via inverses while respecting positional order. \\
\textbf{Correct Solution:} Left-multiply both sides of the equation by $A^{-1}$: $A^{-1}(AX) = A^{-1}B \implies X = A^{-1}B$.

\vspace{5pt}
\noindent
\textbf{29.} \textbf{Flaw:} The student executed the exact opposite of the $2 \times 2$ formula. They swapped the off-diagonal items and negated the main diagonal items. \\
\textbf{Correct Solution:} You must swap the main diagonal items and negate the off-diagonal items. $A^{-1} = \begin{bmatrix} 1 & -2 \\ -3 & 7 \end{bmatrix}$.

\vspace{5pt}
\noindent
\textbf{30.} \textbf{Flaw:} The student incorrectly factored a constant out of the inverse operator. The inverse of a scalar multiplied by a matrix dictates mapping the scalar to its reciprocal. \\
\textbf{Correct Solution:} $M^{-1} = (\frac{1}{2}A)^{-1} = \frac{1}{1/2} A^{-1} = 2 A^{-1}$. Therefore, $M^{-1} = 2 \begin{bmatrix} 4 & 2 \\ 6 & 8 \end{bmatrix} = \begin{bmatrix} 8 & 4 \\ 12 & 16 \end{bmatrix}$.

\newpage

\section{Concept Checklist \& Question Mapping (Appendix)}

For easy reference, below is the Concept Checklist mapped to the relevant question numbers.

\begin{itemize}
    \item \textbf{Concept 1: $2 \times 2$ Inverse Formula Application:} Questions 1, 2, 3, 4, 31
    \item \textbf{Concept 2: Verification and Definition Testing:} Questions 5, 6
    \item \textbf{Concept 3: Solving Linear Systems via Inverses:} Questions 7, 8, 9
    \item \textbf{Concept 4: $n \times n$ Inverse Algorithm:} Questions 10, 11, 12
    \item \textbf{Concept 5: Targeted Inverse Columns:} Questions 13, 14
    \item \textbf{Concept 6: Matrix Equations and Reversals:} Questions 15, 16
    \item \textbf{Concept 7: Algebraic Proofs:} Questions 17, 18, 19, 32
    \item \textbf{Concept 8: Conceptual Synthesis (True/False):} Questions 20, 21, 22, 23, 24
    \item \textbf{Concept 9: Block Matrix Inversion:} Questions 25, 26
    \item \textbf{Concept 10: Find the Flaw (Diagnostic Testing):} Questions 27, 28, 29, 30
\end{itemize}

\end{document}